\documentclass[11pt]{article}
\usepackage{mathblatt}
\begin{document}
\weit
\typeout{PROBE-BEGIN bruchrechnung-e1-k3-s0-v1}
\begin{aufgabe}{\detokenize{bruchrechnung-e1-k3-s0-v1}}
\begin{teile}
\teil Probe

\bruchrechteck{0}{6}
\end{teile}
\end{aufgabe}
\typeout{PROBE-END bruchrechnung-e1-k3-s0-v1}
\typeout{PROBE-BEGIN brueche-dezimalzahlen-e1-k1-s0-v3}
\begin{aufgabe}{\detokenize{brueche-dezimalzahlen-e1-k1-s0-v3}}
\begin{teile}
\teil Probe

\bruchrechteck{0}{5}
\end{teile}
\end{aufgabe}
\typeout{PROBE-END brueche-dezimalzahlen-e1-k1-s0-v3}
\typeout{PROBE-BEGIN prozentrechnung-e1-k3-s3-v3}
\begin{aufgabe}{\detokenize{prozentrechnung-e1-k3-s3-v3}}
\begin{teile}
\teil Probe

\bruchrechteck[5]{0}{20}
\end{teile}
\end{aufgabe}
\typeout{PROBE-END prozentrechnung-e1-k3-s3-v3}
\typeout{PROBE-BEGIN bruchrechnung-e1-k3-s0-v2}
\begin{aufgabe}{\detokenize{bruchrechnung-e1-k3-s0-v2}}
\begin{teile}
\teil Probe

\bruchrechteck{0}{8}
\end{teile}
\end{aufgabe}
\typeout{PROBE-END bruchrechnung-e1-k3-s0-v2}
\typeout{PROBE-BEGIN brueche-dezimalzahlen-e1-k1-s1-v1}
\begin{aufgabe}{\detokenize{brueche-dezimalzahlen-e1-k1-s1-v1}}
\begin{teile}
\teil Probe

\bruchrechteck{3}{5}
\end{teile}
\end{aufgabe}
\typeout{PROBE-END brueche-dezimalzahlen-e1-k1-s1-v1}
\end{document}
